\documentclass[a4paper]{scrartcl}
\usepackage[utf8]{inputenc}
\usepackage{amsmath,amssymb,amsthm}
\usepackage{enumitem,setspace}
\usepackage[usenames,dvipsnames,table]{xcolor}
\usepackage{graphicx}
\DeclareGraphicsExtensions{.pdf,.png,.eps,.jpg}
\usepackage{algorithm2e}

% --- math operators ---
\usepackage{mathtools}
\DeclarePairedDelimiter\set{\{}{\}} % use $\set{1, 2, 3}$
\DeclarePairedDelimiter\abs{\lvert}{\rvert}
\DeclarePairedDelimiter\norm{\lVert}{\rVert}
\DeclarePairedDelimiter\ceils{\lceil}{\rceil}
\DeclarePairedDelimiter\floor{\lfloor}{\rfloor}
\def\then{\ensuremath{\Rightarrow}} % =>
\def\iff{\ensuremath{\Leftrightarrow}} % if and only if <=>
\def\to{\ensuremath{\rightarrow}} % ->
\def\ot{\ensuremath{\leftarrow}} % ->
\def\Oh{\ensuremath{\mathcal{O}}} % big O like $\Oh(n)$
% ---
\DeclareMathOperator{\preorder}{preorder}
\DeclareMathOperator{\postorder}{postorder}
\DeclareMathOperator{\depth}{depth}

% --- FILL OUT THESE 3 FIELDS ---
\def\sheetNumber{1}
\def\names{NAMEN} 
\def\sumPoints{20} 

% --- ### BEGIN DOCUMENT ### ---
\begin{document} 

\begin{doublespace} % header
\noindent \textbf{\LARGE{Übungsblatt \sheetNumber}}
\hfill  {\large Punkte: $\boxed{\qquad  /\; \sumPoints}$}\\
{\Large \textbf{Fortgeschrittene Algorithmen (Wintersemester 2022/23)}\\
Bearbeitet von: \textbf{\names}\\
\today}
\end{doublespace}


\section*{Aufgabe 1: Tamassia-Netzwerk für Rechteck-Knoten}

\subsection*{Idee}

Die Aufgabe ist es, Knoten mit Grad $> 4$ nicht als Punkte, sondern als Rechtecke zu zeichnen. Dafür ersetzen wir einen solchen Knoten $v$ durch einen Ring von $\deg(v)$ Knoten, die auf den vier Seiten eines Rechtecks verteilt werden (nicht in den Ecken!).

Das klassische Tamassia-Netzwerk muss angepasst werden, um sicherzustellen, dass:
\begin{itemize}
    \item Der Ring als Rechteck gezeichnet wird
    \item Jede Seite mindestens einen Knoten hat
    \item Die vier Ecken \textbf{keine} Knoten enthalten
\end{itemize}

\subsection*{Lösung}

Die Anpassung des Tamassia-Netzwerks funktioniert wie folgt:

\textbf{1. Ring-Knoten-Anordnung:}
Die $k = \deg(v)$ Knoten $v_1, \ldots, v_k$ werden auf den vier Seiten angeordnet:
\begin{itemize}
    \item Oben, Rechts, Unten, Links: jeweils $\geq 1$ Knoten
    \item Die Ecken bleiben leer
\end{itemize}

\textbf{2. Winkel-Constraints für Ringknoten:}
Ein Ringknoten $v_i$ auf einer Seite (z.B. oben) hat einen Innenwinkel von $180°$. Das bedeutet, alle seine Kanten liegen auf derselben Seite des Rechtecks:
$$\sum_{e \text{ inzident an } v_i} \text{angle}(e) = 180°$$

\textbf{3. Eck-Verbote (Hauptidee):}
Um zu erzwingen, dass die Ecken keine Knoten erhalten, belegen wir die Kanten, die zu den Ecken führen, mit \textbf{prohibitiv hohen Kosten} oder setzen ihre Kapazität auf 0 im Flussnetzwerk:
$$\text{cost}(\text{Kanten zu Ecken}) = \infty \quad \text{(oder Kapazität } = 0\text{)}$$

Der Min-Cost-Max-Flow-Algorithmus wird dann automatisch diese Kanten vermeiden.

\subsection*{Formalisierung}

Im Tamassia-Netzwerk:
\begin{enumerate}
    \item Für jeden Ringknoten $v_i$ gibt es einen Knoten im Netzwerk
    \item Zwischen Ringknoten und angrenzenden Facetten verlaufen Standard-Tamassia-Kanten
    \item Zusätzlich gibt es spezielle Kanten zu den vier Eck-Positionen mit Kosten $M = \infty$ (oder Kapazität 0)
    \item Der Algorithmus findet automatisch eine knickminimale Zeichnung, bei der die Ecken leer bleiben
\end{enumerate}

\subsection*{Beispiel}

Knoten mit Grad 6 $\to$ Ring $\{v_1, \ldots, v_6\}$:
\begin{itemize}
    \item Oben: $v_1, v_2$
    \item Rechts: $v_3$
    \item Unten: $v_4, v_5$
    \item Links: $v_6$
\end{itemize}
Die Ecken sind unbesetzt. Dadurch entsteht ein Rechteck ohne Knoten in den Ecken.



\section*{Aufgabe 2}

\section*{Aufgabe 3}

\section*{Aufgabe 4}


% -- code for figures
% \begin{figure}[htb]
% \centering
% \includegraphics{filename} % use option [width=\linewidth] to change size to linewidth (or 0.6\linewidth for example)
% \label{fig:name}
% \end{figure}


\end{document}
% --- ### END DOCUMENT ### ---